% entity-id: prop-partial-order
% entity-type: prop
% macro: \PartialOrder
% args: 1
% notation: #1
\hypertarget{prop-partial-order}{}
\begin{theorem}[Partial Order Properties]
\[
\TerForall S \in \Set{\text{Set}}
\TerForall R \subseteq S \CartesianProduct{\times} S
\left(
\SetPartialOrder{\mathrm{IsPartialOrder}}(R, S) \TermEquivalence
\left(
\left( \TerForall a \in S \quad (a, a) \in R \right) \And
\left( \TerForall a, b \in S \quad ((a, b) \in R \And (b, a) \in R) \TermImplication a = b \right) \And
\left( \TerForall a, b, c \in S \quad ((a, b) \in R \And (b, c) \in R) \TermImplication (a, c) \in R \right)
\right)
\right)
\]
\end{theorem}

\begin{proof}[RU]
Отношение \SetPartialOrder{частичного порядка} на \Set{множестве} $S$ — это бинарное отношение $R$, которое одновременно является рефлексивным, антисимметричным и транзитивным.
\end{proof}

\begin{proof}[EN]
A partial order on a set $S$ is a binary relation $R$ that is simultaneously reflexive, antisymmetric, and transitive.
\end{proof}